\documentclass[10pt,a4paper]{amsart}
\usepackage{amsmath,amssymb,amsthm,graphicx}
% Retained preamble for the Stacks Project chapter {CH} (shared project preamble).
% The chapter file pinned-original-source.tex only carries \input{{preamble}} on its first line,
% so no preamble text lives inside the chapter itself.  The Stacks Project uses one single shared
% preamble for every chapter: the snapshot stages it three times (sources/stacks-chow/preamble.tex,
% sources/stacks-etale-cohomology/preamble.tex and focused-seeds/AGP/sources/stacks-snapshot-20261011/
% preamble.tex) and those copies are identical in size.  The block below is copied mechanically,
% line for line, from sources/stacks-chow/preamble.tex lines 14-18 and 154-217: the Xy-pic setup and
% the native theorem declarations together with the mathematical macros.  Nothing was added, removed
% or redefined; the wrapper class and the amsmath/amssymb/amsthm/graphicx base packages are supplied
% by the assembler.
% For commutative diagrams we use Xy-pic
% Chapter: stacks-cohomology.tex -- Cohomology of Sheaves.
\usepackage[all]{xy}

% We use 2cell for 2-commutative diagrams.
\xyoption{2cell}

% Theorem environments.
%
\theoremstyle{plain}
\newtheorem{theorem}[subsection]{Theorem}
\newtheorem{proposition}[subsection]{Proposition}
\newtheorem{lemma}[subsection]{Lemma}

\theoremstyle{definition}
\newtheorem{definition}[subsection]{Definition}
\newtheorem{example}[subsection]{Example}
\newtheorem{exercise}[subsection]{Exercise}
\newtheorem{situation}[subsection]{Situation}

\theoremstyle{remark}
\newtheorem{remark}[subsection]{Remark}
\newtheorem{remarks}[subsection]{Remarks}

\numberwithin{equation}{subsection}

% Macros
%
\def\lim{\mathop{\mathrm{lim}}\nolimits}
\def\colim{\mathop{\mathrm{colim}}\nolimits}
\def\Spec{\mathop{\mathrm{Spec}}}
\def\Hom{\mathop{\mathrm{Hom}}\nolimits}
\def\Ext{\mathop{\mathrm{Ext}}\nolimits}
\def\SheafHom{\mathop{\mathcal{H}\!\mathit{om}}\nolimits}
\def\SheafExt{\mathop{\mathcal{E}\!\mathit{xt}}\nolimits}
\def\Sch{\mathit{Sch}}
\def\Mor{\mathop{\mathrm{Mor}}\nolimits}
\def\Ob{\mathop{\mathrm{Ob}}\nolimits}
\def\Sh{\mathop{\mathit{Sh}}\nolimits}
\def\NL{\mathop{N\!L}\nolimits}
\def\CH{\mathop{\mathrm{CH}}\nolimits}
\def\proetale{{pro\text{-}\acute{e}tale}}
\def\etale{{\acute{e}tale}}
\def\QCoh{\mathit{QCoh}}
\def\Ker{\mathop{\mathrm{Ker}}}
\def\Im{\mathop{\mathrm{Im}}}
\def\Coker{\mathop{\mathrm{Coker}}}
\def\Coim{\mathop{\mathrm{Coim}}}

% Boxtimes
%
\DeclareMathSymbol{\boxtimes}{\mathbin}{AMSa}{"02}

%
% Macros for moduli stacks/spaces
%
\def\QCohstack{\mathcal{QC}\!\mathit{oh}}
\def\Cohstack{\mathcal{C}\!\mathit{oh}}
\def\Spacesstack{\mathcal{S}\!\mathit{paces}}
\def\Quotfunctor{\mathrm{Quot}}
\def\Hilbfunctor{\mathrm{Hilb}}
\def\Curvesstack{\mathcal{C}\!\mathit{urves}}
\def\Polarizedstack{\mathcal{P}\!\mathit{olarized}}
\def\Complexesstack{\mathcal{C}\!\mathit{omplexes}}
% \Pic is the operator that assigns to X its picard group, usage \Pic(X)
% \Picardstack_{X/B} denotes the Picard stack of X over B
% \Picardfunctor_{X/B} denotes the Picard functor of X over B
\def\Pic{\mathop{\mathrm{Pic}}\nolimits}
\def\Picardstack{\mathcal{P}\!\mathit{ic}}
\def\Picardfunctor{\mathrm{Pic}}
\def\Deformationcategory{\mathcal{D}\!\mathit{ef}}
\usepackage{cleveref}
\crefname{originalequation}{original equation}{original equations}
\Crefname{originalequation}{Original equation}{Original equations}
\makeatletter
\newlabel{equation-projection-formula-map}{{0B53}{source}{}{sourcepointer-1}{}}
\newlabel{equation-projection-formula-map@cref}{{[originalequation][0][]0B53}{[1][0][]source}{}{}{}}
\makeatother
\crefname{originalreference}{original reference}{original references}
\Crefname{originalreference}{Original reference}{Original references}
\makeatletter
\newlabel{modules-lemma-i-star-exact}{{01AX}{source}{}{sourcepointer-2}{}}
\newlabel{modules-lemma-i-star-exact@cref}{{[originalreference][0][]01AX}{[1][0][]source}{}{}{}}
\makeatother
\crefname{originallemma}{original lemma}{original lemmas}
\Crefname{originallemma}{Original lemma}{Original lemmas}
\makeatletter
\newlabel{lemma-pullback-K-flat}{{06YC}{source}{}{sourcepointer-3}{}}
\newlabel{lemma-pullback-K-flat@cref}{{[originallemma][0][]06YC}{[1][0][]source}{}{}{}}
\makeatother
\crefname{originalreference}{original reference}{original references}
\Crefname{originalreference}{Original reference}{Original references}
\makeatletter
\newlabel{modules-lemma-i-star-right-adjoint}{{01AZ}{source}{}{sourcepointer-4}{}}
\newlabel{modules-lemma-i-star-right-adjoint@cref}{{[originalreference][0][]01AZ}{[1][0][]source}{}{}{}}
\makeatother
\crefname{originalreference}{original reference}{original references}
\Crefname{originalreference}{Original reference}{Original references}
\makeatletter
\newlabel{sheaves-lemma-stalks-closed-pushforward}{{00AE}{source}{}{sourcepointer-5}{}}
\newlabel{sheaves-lemma-stalks-closed-pushforward@cref}{{[originalreference][0][]00AE}{[1][0][]source}{}{}{}}
\makeatother
\crefname{originalreference}{original reference}{original references}
\Crefname{originalreference}{Original reference}{Original references}
\makeatletter
\newlabel{sheaves-lemma-stalk-pullback-modules}{{0098}{source}{}{sourcepointer-6}{}}
\newlabel{sheaves-lemma-stalk-pullback-modules@cref}{{[originalreference][0][]0098}{[1][0][]source}{}{}{}}
\makeatother
\title{The projection formula map is an isomorphism for a morphism of ringed spaces that is a homeomorphism onto a closed subset}
\author{The Stacks Project Authors}
\date{stacks-project:cohomology.tex:local-snapshot-20261011; GNU FDL 1.2 or later}
\begin{document}
\maketitle
\noindent\textit{Editorial note.} All delivered mathematical text is the author's own, copied verbatim from the pinned Stacks Project chapter file cohomology.tex; no statement or proof step was rewritten, reordered or removed.\par\medskip
\noindent\textit{Editorial note.} The counted span cites six cross-chapter labels, including the equation defining the projection formula map; each is recorded in assembly.external\_references with its canonical tag as a navigation pointer.\par\medskip
\noindent\textit{Editorial note.} The proof cites no bibliography item inside the counted span.\par\medskip
\noindent\textit{Allowed standard prerequisites (editorial).}\begin{itemize}
\item Exactness of pushforward along a homeomorphism onto a closed subset. Stacks Project, Modules on Sites, Lemma modules-lemma-i-star-exact (tag 01AX). Allows replacing Rf\_* by f\_* on any representative complex, which is what makes the comparison a levelwise computation.
\item K-flatness and the projection formula for pullback. Stacks Project, Cohomology of Sheaves, Lemma lemma-pullback-K-flat (tag 06YC). Supplies a representative of Lf\textasciicircum{}*K with which the tensor product computation can be performed on the level of complexes.
\item Stalks of closed pushforward and of pullback of modules. Stacks Project, Sheaves, Lemmas sheaves-lemma-stalks-closed-pushforward (tag 00AE) and sheaves-lemma-stalk-pullback-modules (tag 0098). Identify the stalks appearing in the final computation of the single-module identity.
\end{itemize}\medskip
\ifdefined\hypertarget\hypertarget{sourcepointer-1}{}\fi
\noindent\textit{Source reference.} Original equation 0B53: sources/stacks-chow/tags.txt, source lines 14021--14021. Cross-chapter standard prerequisite cited by the author inside the counted span; it is used as a pointer and its own proof is not part of this item.\par\medskip
\ifdefined\hypertarget\hypertarget{sourcepointer-2}{}\fi
\noindent\textit{Source reference.} Original reference 01AX: sources/stacks-chow/tags.txt, source lines 1706--1706. Cross-chapter standard prerequisite cited by the author inside the counted span; it is used as a pointer and its own proof is not part of this item.\par\medskip
\ifdefined\hypertarget\hypertarget{sourcepointer-3}{}\fi
\noindent\textit{Source reference.} Original lemma 06YC: sources/stacks-chow/tags.txt, source lines 8881--8881. Cross-chapter standard prerequisite cited by the author inside the counted span; it is used as a pointer and its own proof is not part of this item.\par\medskip
\ifdefined\hypertarget\hypertarget{sourcepointer-4}{}\fi
\noindent\textit{Source reference.} Original reference 01AZ: sources/stacks-chow/tags.txt, source lines 1708--1708. Cross-chapter standard prerequisite cited by the author inside the counted span; it is used as a pointer and its own proof is not part of this item.\par\medskip
\ifdefined\hypertarget\hypertarget{sourcepointer-5}{}\fi
\noindent\textit{Source reference.} Original reference 00AE: sources/stacks-chow/tags.txt, source lines 386--386. Cross-chapter standard prerequisite cited by the author inside the counted span; it is used as a pointer and its own proof is not part of this item.\par\medskip
\ifdefined\hypertarget\hypertarget{sourcepointer-6}{}\fi
\noindent\textit{Source reference.} Original reference 0098: sources/stacks-chow/tags.txt, source lines 344--344. Cross-chapter standard prerequisite cited by the author inside the counted span; it is used as a pointer and its own proof is not part of this item.\par\medskip
\begin{lemma}
\label{lemma-projection-formula-closed-immersion}
Let $f : X \to Y$ be a morphism of ringed spaces such that $f$ is a
homeomorphism onto a closed subset. Then
(\ref{equation-projection-formula-map}) is an isomorphism always.
\end{lemma}

\begin{proof}
Since $f$ is a homeomorphism onto a closed subset, the functor $f_*$
is exact (Modules, Lemma \ref{modules-lemma-i-star-exact}). Hence
$Rf_*$ is computed by applying $f_*$ to any representative complex.
Choose a K-flat complex $\mathcal{K}^\bullet$ of $\mathcal{O}_Y$-modules
representing $K$ and choose any complex $\mathcal{E}^\bullet$
of $\mathcal{O}_X$-modules representing $E$. Then
$Lf^*K$ is represented by $f^*\mathcal{K}^\bullet$ which is
a K-flat complex of $\mathcal{O}_X$-modules
(Lemma \ref{lemma-pullback-K-flat}). Thus the right hand side of
(\ref{equation-projection-formula-map}) is represented by
$$
f_*\text{Tot}(\mathcal{E}^\bullet
\otimes_{\mathcal{O}_X} f^*\mathcal{K}^\bullet)
$$
By the same reasoning we see that the left hand side is represented by
$$
\text{Tot}(f_*\mathcal{E}^\bullet \otimes_{\mathcal{O}_Y} \mathcal{K}^\bullet)
$$
Since $f_*$ commutes with direct sums
(Modules, Lemma \ref{modules-lemma-i-star-right-adjoint})
it suffices to show that
$$
f_*(\mathcal{E} \otimes_{\mathcal{O}_X} f^*\mathcal{K}) =
f_*\mathcal{E} \otimes_{\mathcal{O}_Y} \mathcal{K}
$$
for any $\mathcal{O}_X$-module $\mathcal{E}$ and $\mathcal{O}_Y$-module
$\mathcal{K}$. We will check this by checking on stalks.
Let $y \in Y$. If $y \not \in f(X)$, then the stalks
of both sides are zero. If $y = f(x)$, then we see that we have to show
$$
\mathcal{E}_x \otimes_{\mathcal{O}_{X, x}}
(\mathcal{O}_{X, x} \otimes_{\mathcal{O}_{Y, y}} \mathcal{F}_y) =
\mathcal{E}_x \otimes_{\mathcal{O}_{Y, y}} \mathcal{F}_y
$$
(using Sheaves, Lemma \ref{sheaves-lemma-stalks-closed-pushforward}
and Lemma \ref{sheaves-lemma-stalk-pullback-modules}).
This equality holds and therefore the lemma has been proved.
\end{proof}

\end{document}
